\documentclass{subfiles}
\begin{document}
    Later in these notes, see \Cref{sec:EllipticCurvesAlgos},
        we will consider some algorithms that are based on the notion of projective space.
        Hence, for the sake of completeness, 
        in this section we provide the reader with the fundamentals of this field.

    Although it has a deep foundation in math,
        the notion of projective space was first used in art to represent objects that are far from the viewer.
        In simpler terms, a projective space is what we usually refer to as perspective.

    More formally, a projective space over a field $\K$ is the set 
        $$
            \ps^{n} = \set[\vec{0} \ne \vec{v} \in \K^{n + 1}]{P = \gen{\vec{v}}} 
        $$
        where $\gen{\vec{v}} = \set[\lambda \in \K]{\lambda \vec{v}}$.
        Let $P$ be a point in $\ps^{n}$.
        We say that $P$ is in homogeneous coordinates, if its components $x_{0}, \ldots, x_{n}$ differ,
        at most, by a proportionality factor. Viceversa, we sau that $P$ is given in non-homogeneous coordinates.
        In the latter case $P$ can assume two forms:
        $$
            P = \begin{cases}
                [1, \frac{x_{1}}{x_{0}}, \ldots, \frac{x_{n}}{x_{0}}], & \text{if } x_{0} \ne 0; \\ 
                [0, x_{1}, \ldots, x_{n}], & \text{otherwise}.
            \end{cases}
        $$
        To a point $P$ of the form $[1, \sfrac{x_{1}}{x_{0}}, \ldots, \sfrac{x_{n}}{n}]$,
        we can associate a point $P(X, \ldots, Z)$ where $X = \sfrac{x_{1}}{x_{0}}, \cdots, Z = \sfrac{x_{n}}{x_{0}}$;
        we call these points \emph{affine} (or proper) points. All non-affine points are called points at infinity. 

    Now, let us consider what happens with polynomial in the projective space.
        By the way $\ps^{n}$ is defined, we must have that given two point $P, Q$,
        with $P$ and $Q$ proportional, the evaluation of a polynomial at $P$ and $Q$ gives the same
        result, since $P$ and $Q$ represent the same point.
        As it turns out, this is not the case unless the polynomial/function is homogeneous;
        that is, each of the monomyals that compose it, have the same degree.


    % TODO: section on elliptic curves.
        
\end{document}
